\documentclass[10pt,a4paper]{amsart}
\usepackage[margin=19mm]{geometry}
\usepackage{amsmath,amssymb,mathtools}
\usepackage[scr=rsfs,cal=euler]{mathalfa}
\usepackage{xfrac}
\usepackage[unicode,hidelinks,bookmarks=false]{hyperref}
\newcommand{\ie}{i.e.\ }
\newcommand{\cf}{cf.\ }
\newcommand{\resp}{respectively\ }
\newcommand{\Subsection}{\subsection}
\providecommand{\nobreakdash}{\nobreak\hbox{-}}
\setlength{\emergencystretch}{2em}
\multlinegap0pt
\usepackage{bm}
\usepackage{bbm}
\usepackage{dsfont}
\usepackage{xspace}
\usepackage{cancel}
\usepackage{mathtools}
\usepackage{amsthm}
\makeatletter
\@ifundefined{theorem}{\newtheorem{theorem}{Theorem}}{}
\makeatother
\renewcommand{\hat}{\widehat}
\title[Large deviations and the matrix product ansatz]{Large deviations and the matrix product ansatz}
\author{Davide Gabrielli, Federica Iacovissi}
\date{Source: arXiv:2601.11820v1 (CC BY 4.0)}
\begin{document}
\maketitle

\noindent\emph{Source and licence.} Large deviations and the matrix product ansatz. Version arXiv:2601.11820v1 (CC BY 4.0), distributed under the Creative Commons Attribution 4.0 International licence, as recorded in the arXiv licence metadata for this exact version and shown on the abstract page https://arxiv.org/abs/2601.11820v1. Full rights notice: licenses/arxiv-2601.11820-CC-BY-4.0.txt.\par\smallskip

\noindent\textit{Editorial scope.} Counted unit: the Theorem whose source label is \texttt{None} in the article (source0.tex lines 1259--1270, source label \texttt{None}), delivered together with the article's own complete original proof of it (source0.tex lines 1271--1295, one \texttt{proof} environment of 25 lines). No mathematics is written, repaired or re-derived: statement and proof are the article's own text line for line. Required context retained uncounted: imp (lines 932--934); rhoezeta (lines 1229--1231). Every definition, display and label of those lines is retained unchanged. Omitted from this package and not redistributed: 1--931, 935--1228, 1232--1258, 1296--1512. Nothing inside a retained excerpt is omitted or altered: every retained body is delivered exactly as the article's lines are written, so no omission receipt of any kind accompanies this package. Citations: the retained excerpts carry no citation command at all, so no bibliography entry of the article is transcribed and no reference is redistributed. Reference adaptation: none was needed and none was made; every reference inside the delivered unit resolves inside it, so no unresolved reference or citation is produced. Printed slips kept literal: no printed word, symbol, hypothesis or number of the counted statement or of its proof was altered. Layout and class: the article's own class is \texttt{amsart} with packages including amsmath, amssymb, cite, euscript, mathtools, xcolor, braket, stackrel, tikz, ulem; the wrapper is the lane's pinned \texttt{amsart} editorial wrapper at 10pt on A4, which restates the theorem-like environments the retained text needs and adds the article's own macro definitions that the retained text needs (\texttt{\textbackslash hat}), with extra wrapper packages (bm, bbm, dsfont, xspace, cancel, mathtools). The delivered numbering is the unit's own. Assets: no retained span contains a figure environment, a graphics-inclusion command, a PDF-page inclusion, a table or a TikZ picture; no image file is copied into this package, the package declares no asset dependency and no image file is redistributed with it. Licence: version arXiv:2601.11820v1 of this article is distributed under the Creative Commons Attribution 4.0 International licence, as recorded in the arXiv licence metadata for this exact version and shown on the abstract page https://arxiv.org/abs/2601.11820v1. A full rights notice is delivered at licenses/arxiv-2601.11820-CC-BY-4.0.txt. Characters: every retained excerpt is pure ASCII, so the delivered engine, XeLaTeX with its default Latin Modern text font, reports no missing character.
	\begin{equation}\label{imp}
		\hat \pi_N[\eta]=\sum_{a}\sum_{ b}\hat \Pi_N[\eta,\zeta]\left(a, b\right)a\,;
	\end{equation}

		\begin{equation}\label{rhoezeta}
			I^{\mathfrak S}(z,\rho)=z(0)\log\frac{\alpha}{1-\alpha}+z(1)\log\frac{\beta}{1-\beta}+\int_0^1 \Big[H\big(\rho(x)\big)+H\Big(\rho(x)-\dot{z}(x)\Big)\Big] \, dx +  C'
		\end{equation}

	\begin{theorem}
		Let $F(x)=\int_0^x\rho(y)dy$ and $\mathcal G$ the set of absolutely continuous functions $G$ such that $G(0)=0$ and $0\leq \dot G(x)\leq 1$. Let $\eta$ be distributed according to the invariant measure of the boundary driven TASEP when the boundary parameters are such that $ab<1$. We have that the large deviations rate functional for $\hat\pi_N$ in \eqref{imp} is given by
		\begin{equation*}
			\begin{aligned}
				I^{\mathrm{TASEP}}(\rho)
				= \inf_{G\in \mathcal G} \biggl\{
				&\int_0^1 \Bigl[H\left(\dot F(y)\right) + H\left(\dot{G}(y)\right)\Bigr]\, dy  \\
				&\quad + \ln(ab)\, \min_{x\in[0,1]} \bigl[F(x)-G(x)\bigr] - \ln(b)\, \bigl(F(1)-G(1)\bigr)
				\biggr\}.
			\end{aligned}
		\end{equation*}
	\end{theorem}

	\begin{proof}
		By the contraction principle  we have to  minimize the functional \eqref{rhoezeta} over the functions $z(x) \ge 0$ such that $0 \le \rho(x)-\dot z(x) \le 1$ for almost any $x\in[0,1]$.
		
		We define:
		\begin{align*}
			& F(x)= \int_0^x \rho(y)\,dy \\
			&G(x)=F(x)-\Big(z(x)-z(0)\Big)\,,
		\end{align*}
		and note that $G(0)=0$ and $\dot G(x)=\rho(x)-\dot z(x)$ so that $G\in \mathcal G$; moreover $F(x)+z(0)-G(x)=z(x)\geq 0$.
		The functional \eqref{rhoezeta} can be written in terms of $F, G, z(0)$ as
		$$
		\int_0^1 \Bigl[H\left(\dot F(y)\right) + H\left(\dot{G}(y)\right)\Bigr]\, dy  
		-z(0)\ln(ab)-\bigl[F(1)-G(1)\bigr]\ln(b):=\mathcal L(F, G, z(0))\,.
		$$
		By the contraction principle we have
		We rewrite the problem:
		\begin{equation*}
			I^{\mathrm{TASEP}}(\rho)=\inf_{\big\{z(0) \ge 0\big\}}  \inf_{\big\{G\in \mathcal G : G(x) \le F(x)+z(0)\big\}} \mathcal{L}(F, G, z(0)) \,.
		\end{equation*}
		We exchange the order of the infima obtaining that the above variational problem is equivalent to
		\begin{equation*}
			I^{\mathrm{TASEP}}(\rho)=\inf_{\big\{G\in \mathcal G\big\}}  \inf_{\big\{z(0)\ge \max_{x \in [0,1]}[G(x)-F(x)]\big\}} \mathcal{L}(F, G, z(0)) \,.
		\end{equation*}
		Since $ab<1$ the value of $z(0)$ that minimizes the first infimum is exactly $z(0)= \max_{x \in [0,1]}[G(x)-F(x)]$ and this finishes the proof.
	\end{proof}

\end{document}
